% BP_RECORD_BEGIN
% {
%   "id": "weak_comparison",
%   "title": "Weak Comparison",
%   "order": 2,
%   "summary": "Generic C^3 Lindeberg comparison and its project-kernel specialization for psi versus gaussianPsi.",
%   "reader_notes": "This chapter isolates the invariance-principle input. It shows both the paper lemma and the project-specific specialization that turns the generic theorem into the psi-versus-gaussian comparison used later.",
%   "record_type": "chapter",
%   "chapter_file": "blueprint/src/chapters/weak_comparison.tex"
% }
% BP_RECORD_END
\section{Weak Comparison}
\label{chap:weak-comparison}

This chapter records the generic weak-comparison theorem and the specialization
route that produces the project-level comparison for
$\psi_n - \psi_{G,n}$. The blueprint keeps the generic theorem and the project
specialization separate because the Lean development really uses both.

% BP_RECORD_BEGIN
% {
%   "id": "paper-lem-weak-comparison",
%   "label": "lem:weak-comparison",
%   "title": "Lindeberg comparison",
%   "kind": "lemma",
%   "tex_path": "documents/short.tex",
%   "statement_lines": {
%     "start": 1436,
%     "end": 1448
%   },
%   "proof_lines": {
%     "start": 1449,
%     "end": 1492
%   },
%   "lean_primary": "lean-weak-c3",
%   "lean_support": [
%     "lean-weak-moo"
%   ],
%   "module_owner": "module-weak-invariance",
%   "summary": "The generic weak-comparison theorem formalizes the paper\u2019s C^3 Lindeberg replacement bound for quadratic forms.",
%   "proof_steps": [
%     {
%       "id": "step-weak-decomposition",
%       "title": "Decomposition Q_f = U + z_m V",
%       "tex_lines": {
%         "start": 1452,
%         "end": 1459
%       },
%       "lean_refs": [
%         "lean-weak-c3"
%       ],
%       "summary": "Split the quadratic form into a part independent of the replaced coordinate and a linear form in that coordinate."
%     },
%     {
%       "id": "step-weak-taylor",
%       "title": "Third-order Taylor step",
%       "tex_lines": {
%         "start": 1460,
%         "end": 1468
%       },
%       "lean_refs": [
%         "lean-weak-c3"
%       ],
%       "summary": "Apply third-order Taylor expansion and use matching first two moments of the sign and Gaussian variables."
%     },
%     {
%       "id": "step-weak-adjacent",
%       "title": "Adjacent replacement estimate",
%       "tex_lines": {
%         "start": 1469,
%         "end": 1490
%       },
%       "lean_refs": [
%         "lean-weak-c3"
%       ],
%       "summary": "Control each replacement step by a cubic influence term via second and fourth moment estimates for the linear form."
%     },
%     {
%       "id": "step-weak-telescope",
%       "title": "Telescoping over m",
%       "tex_lines": {
%         "start": 1449,
%         "end": 1451
%       },
%       "lean_refs": [
%         "lean-weak-c3"
%       ],
%       "summary": "Write the full difference as a telescoping sum of adjacent replacement steps."
%     }
%   ],
%   "record_type": "paper_item",
%   "chapter": "weak_comparison"
% }
% BP_RECORD_END
\begin{lemma}[Lindeberg comparison]
\label{lem:weak-comparison}
\lean{quadraticForm_lindeberg_comparison_C3}
\leanok
\uses{lean-weak-moo}
The generic weak-comparison theorem formalizes the paper’s C^3 Lindeberg replacement bound for quadratic forms.
\end{lemma}

% BP_RECORD_BEGIN
% {
%   "id": "paper-cor-weak-comparison",
%   "label": "cor:weak-comparison",
%   "title": "Gaussian comparison for \u03c8",
%   "kind": "corollary",
%   "tex_path": "documents/short.tex",
%   "statement_lines": {
%     "start": 1496,
%     "end": 1502
%   },
%   "proof_lines": {
%     "start": 1504,
%     "end": 1515
%   },
%   "lean_primary": "lean-psi-vs-gaussian",
%   "lean_support": [
%     "lean-weak-c3",
%     "lean-weak-moo",
%     "lean-paper-three-half"
%   ],
%   "module_owner": "module-weak-invariance",
%   "summary": "The characteristic-function comparison is derived by specializing the generic C^3 theorem to cos and sin and then rewriting the influence normalization.",
%   "proof_steps": [
%     {
%       "id": "step-psi-cos",
%       "title": "Specialize to cos",
%       "tex_lines": {
%         "start": 1508,
%         "end": 1510
%       },
%       "lean_refs": [
%         "lean-weak-c3",
%         "lean-weak-moo"
%       ],
%       "summary": "Apply the generic comparison theorem to the cosine part of the characteristic function."
%     },
%     {
%       "id": "step-psi-sin",
%       "title": "Specialize to sin",
%       "tex_lines": {
%         "start": 1508,
%         "end": 1510
%       },
%       "lean_refs": [
%         "lean-weak-c3",
%         "lean-weak-moo"
%       ],
%       "summary": "Apply the same theorem to the sine part and combine real and imaginary contributions."
%     },
%     {
%       "id": "step-psi-normalization",
%       "title": "Convert influence normalization",
%       "tex_lines": {
%         "start": 1511,
%         "end": 1514
%       },
%       "lean_refs": [
%         "lean-paper-three-half",
%         "lean-psi-vs-gaussian"
%       ],
%       "summary": "Rewrite the kernel-influence sum in the project\u2019s row-influence normalization."
%     }
%   ],
%   "record_type": "paper_item",
%   "chapter": "weak_comparison"
% }
% BP_RECORD_END
\begin{corollary}[Gaussian comparison for ψ]
\label{cor:weak-comparison}
\lean{psi_sub_gaussianPsi_le_threeHalfInfluenceSum}
\leanok
\uses{lean-weak-c3,lean-weak-moo,lean-paper-three-half}
The characteristic-function comparison is derived by specializing the generic C^3 theorem to cos and sin and then rewriting the influence normalization.
\end{corollary}

% BP_RECORD_BEGIN
% {
%   "id": "lean-row-influence",
%   "symbol": "rowInfluence",
%   "file": "RequestProject/HadamardCn3MOO.lean",
%   "line": 11413,
%   "module": "module-moo",
%   "kind": "definition",
%   "paper_item": null,
%   "used_by": [
%     "lean-three-half-sum"
%   ],
%   "depends_on": [],
%   "summary": "Row influence I_k(lambda) in the project normalization.",
%   "record_type": "lean_item",
%   "chapter": "weak_comparison"
% }
% BP_RECORD_END
\begin{definition}[rowInfluence]
\label{lean-row-influence}
\lean{rowInfluence}
\moduleowner{module-moo}
Row influence I_k(lambda) in the project normalization.
\end{definition}
\leanitemdetails{lean-row-influence}

% BP_RECORD_BEGIN
% {
%   "id": "lean-three-half-sum",
%   "symbol": "threeHalfInfluenceSum",
%   "file": "RequestProject/HadamardCn3MOO.lean",
%   "line": 11417,
%   "module": "module-moo",
%   "kind": "definition",
%   "paper_item": null,
%   "used_by": [
%     "lean-paper-three-half",
%     "lean-weak-moo",
%     "lean-psi-vs-gaussian"
%   ],
%   "depends_on": [
%     "lean-row-influence"
%   ],
%   "summary": "Project statistic \u03a3 I_k(lambda)^(3/2).",
%   "record_type": "lean_item",
%   "chapter": "weak_comparison"
% }
% BP_RECORD_END
\begin{definition}[threeHalfInfluenceSum]
\label{lean-three-half-sum}
\lean{threeHalfInfluenceSum}
\uses{lean-row-influence}
\moduleowner{module-moo}
Project statistic Σ I_k(lambda)^(3/2).
\end{definition}
\leanitemdetails{lean-three-half-sum}

% BP_RECORD_BEGIN
% {
%   "id": "lean-kernel-influence",
%   "symbol": "kernelInfluence",
%   "file": "RequestProject/HadamardCn3MOO.lean",
%   "line": 11514,
%   "module": "module-moo",
%   "kind": "definition",
%   "paper_item": null,
%   "used_by": [
%     "lean-weak-c3",
%     "lean-paper-three-half"
%   ],
%   "depends_on": [],
%   "summary": "Kernel influence used in the generic quadratic-form comparison theorem.",
%   "record_type": "lean_item",
%   "chapter": "weak_comparison"
% }
% BP_RECORD_END
\begin{definition}[kernelInfluence]
\label{lean-kernel-influence}
\lean{kernelInfluence}
\moduleowner{module-moo}
Kernel influence used in the generic quadratic-form comparison theorem.
\end{definition}
\leanitemdetails{lean-kernel-influence}

% BP_RECORD_BEGIN
% {
%   "id": "lean-moo-kernel",
%   "symbol": "mooKernel",
%   "file": "RequestProject/HadamardCn3MOO.lean",
%   "line": 11524,
%   "module": "module-moo",
%   "kind": "definition",
%   "paper_item": null,
%   "used_by": [
%     "lean-weak-moo",
%     "lean-psi-vs-gaussian",
%     "lean-paper-three-half"
%   ],
%   "depends_on": [],
%   "summary": "Project kernel attached to \u03bb in the weak-comparison route.",
%   "record_type": "lean_item",
%   "chapter": "weak_comparison"
% }
% BP_RECORD_END
\begin{definition}[mooKernel]
\label{lean-moo-kernel}
\lean{mooKernel}
\moduleowner{module-moo}
Project kernel attached to λ in the weak-comparison route.
\end{definition}
\leanitemdetails{lean-moo-kernel}

% BP_RECORD_BEGIN
% {
%   "id": "lean-paper-three-half",
%   "symbol": "paper_threeHalfInfluenceSum_eq_eight_kernelInfluence_sum",
%   "file": "RequestProject/HadamardCn3PaperSpec.lean",
%   "line": 48,
%   "module": "module-paper-spec",
%   "kind": "normalization",
%   "paper_item": "paper-cor-weak-comparison",
%   "used_by": [
%     "lean-psi-vs-gaussian"
%   ],
%   "depends_on": [
%     "lean-three-half-sum",
%     "lean-kernel-influence",
%     "lean-moo-kernel"
%   ],
%   "summary": "Specification-layer theorem rewriting the project influence statistic into the paper kernel-influence normalization.",
%   "record_type": "lean_item",
%   "chapter": "weak_comparison"
% }
% BP_RECORD_END
\begin{lemma}[paper_threeHalfInfluenceSum_eq_eight_kernelInfluence_sum]
\label{lean-paper-three-half}
\lean{paper_threeHalfInfluenceSum_eq_eight_kernelInfluence_sum}
\uses{lean-three-half-sum,lean-kernel-influence,lean-moo-kernel}
\moduleowner{module-paper-spec}
Specification-layer theorem rewriting the project influence statistic into the paper kernel-influence normalization.
\end{lemma}
\leanitemdetails{lean-paper-three-half}

% BP_RECORD_BEGIN
% {
%   "id": "lean-weak-c3",
%   "symbol": "quadraticForm_lindeberg_comparison_C3",
%   "file": "RequestProject/HadamardCn3WeakInvariance.lean",
%   "line": 6454,
%   "module": "module-weak-invariance",
%   "kind": "paper_endpoint",
%   "paper_item": "paper-lem-weak-comparison",
%   "used_by": [
%     "lean-weak-moo",
%     "lean-psi-vs-gaussian"
%   ],
%   "depends_on": [
%     "lean-kernel-influence"
%   ],
%   "summary": "Generic C^3 quadratic-form comparison theorem in the paper normalization.",
%   "record_type": "lean_item",
%   "chapter": "weak_comparison"
% }
% BP_RECORD_END
\begin{lemma}[quadraticForm_lindeberg_comparison_C3]
\label{lean-weak-c3}
\lean{quadraticForm_lindeberg_comparison_C3}
\uses{lean-kernel-influence}
\moduleowner{module-weak-invariance}
Generic C^3 quadratic-form comparison theorem in the paper normalization.
\end{lemma}
\leanitemdetails{lean-weak-c3}

% BP_RECORD_BEGIN
% {
%   "id": "lean-weak-moo",
%   "symbol": "quadraticForm_lindeberg_comparison_mooKernel",
%   "file": "RequestProject/HadamardCn3WeakInvariance.lean",
%   "line": 6634,
%   "module": "module-weak-invariance",
%   "kind": "paper_support",
%   "paper_item": "paper-cor-weak-comparison",
%   "used_by": [
%     "lean-psi-vs-gaussian"
%   ],
%   "depends_on": [
%     "lean-weak-c3",
%     "lean-moo-kernel",
%     "lean-three-half-sum"
%   ],
%   "summary": "Project-kernel specialization of the generic weak-comparison theorem.",
%   "record_type": "lean_item",
%   "chapter": "weak_comparison"
% }
% BP_RECORD_END
\begin{lemma}[quadraticForm_lindeberg_comparison_mooKernel]
\label{lean-weak-moo}
\lean{quadraticForm_lindeberg_comparison_mooKernel}
\uses{lean-weak-c3,lean-moo-kernel,lean-three-half-sum}
\moduleowner{module-weak-invariance}
Project-kernel specialization of the generic weak-comparison theorem.
\end{lemma}
\leanitemdetails{lean-weak-moo}

% BP_RECORD_BEGIN
% {
%   "id": "lean-psi-vs-gaussian",
%   "symbol": "psi_sub_gaussianPsi_le_threeHalfInfluenceSum",
%   "file": "RequestProject/HadamardCn3WeakInvariance.lean",
%   "line": 6726,
%   "module": "module-weak-invariance",
%   "kind": "paper_endpoint",
%   "paper_item": "paper-cor-weak-comparison",
%   "used_by": [
%     "lean-residual-far-shell",
%     "lean-residual-estimate"
%   ],
%   "depends_on": [
%     "lean-weak-c3",
%     "lean-weak-moo",
%     "lean-paper-three-half"
%   ],
%   "summary": "Paper-facing Gaussian comparison bound for the characteristic function \u03c8.",
%   "record_type": "lean_item",
%   "chapter": "weak_comparison"
% }
% BP_RECORD_END
\begin{corollary}[psi_sub_gaussianPsi_le_threeHalfInfluenceSum]
\label{lean-psi-vs-gaussian}
\lean{psi_sub_gaussianPsi_le_threeHalfInfluenceSum}
\uses{lean-weak-c3,lean-weak-moo,lean-paper-three-half}
\moduleowner{module-weak-invariance}
Paper-facing Gaussian comparison bound for the characteristic function ψ.
\end{corollary}
\leanitemdetails{lean-psi-vs-gaussian}
